% Provenance (never printed): NSI terminale "Structures de donnees: arbres, arbres binaires, arbres binaires de
% recherche; algorithmique: parcours d'arbres, recherche et insertion dans un ABR". See sources/book2/BRIEFS.md.
\chapter{Trees and Binary Search Trees}\label{ch:g12:trees}

``I am thinking of a number between $1$ and $100$.'' The best first guess is $50$; after ``smaller'', $25$; after
``larger'', $37$. Every game follows one path through a fixed set of questions: $50$ at the top, $25$ and $75$ below
it, then $12$, $37$, $62$, $87$, and so on. Drawn out, these questions form a \omterm{def:g12:trees:tree}{tree}, and the length of the longest
path is the number of guesses the game can need.

This chapter names the parts of a \omterm{def:g12:trees:tree}{tree}, builds \omterm{def:g12:trees:binary-tree}{binary trees} with a class, visits their \omterm{def:g12:trees:tree}{nodes} in four orders, and
stores keys in \omterm{def:g12:trees:bst}{binary search trees}, where finding a key costs as many steps as the \omterm{def:g12:trees:tree}{tree} is tall. Then it asks how
tall such a \omterm{def:g12:trees:tree}{tree} grows, and measures it.

\begin{recall}
A \omterm{def:g12:recursion:recursion}{recursive function} calls itself on smaller inputs, and is proved by induction (chapter \ref{ch:g12:recursion}). A
queue gives back its elements in the order they came (chapter \ref{ch:g12:lists-stacks-queues}). Folders inside
folders (Book 1, grade 3) and the nested elements of a web page (Book 1, grade 9) are \omterm{def:g12:trees:tree}{trees}.
\end{recall}

\section{Vocabulary}\label{sec:g12:trees:vocabulary}

\begin{definition}[Tree, node, root, leaf, parent, child, subtree]\label{def:g12:trees:tree}
A \emph{tree}\index{tree} is a set of \emph{nodes}\index{node} linked so that one node, the
\emph{root}\index{root}, has no parent, and every other node has exactly one \emph{parent}\index{parent}; a node is a
\emph{child}\index{child} of its parent. A node with no child is a \emph{leaf}\index{leaf}. A node together with all
the nodes below it forms a \emph{subtree}\index{subtree}.
\end{definition}

\begin{definition}[Depth, height, size]\label{def:g12:trees:height}
The \emph{depth}\index{depth} of a \omterm{def:g12:trees:tree}{node} is the number of links from the \omterm{def:g12:trees:tree}{root} to it: the \omterm{def:g12:trees:tree}{root} has depth $0$. The
\emph{height}\index{height} of a \omterm{def:g12:trees:tree}{tree} is the largest depth of its \omterm{def:g12:trees:tree}{nodes}; a \omterm{def:g12:trees:tree}{tree} of one \omterm{def:g12:trees:tree}{node} has height $0$, and by
convention the empty \omterm{def:g12:trees:tree}{tree} has height $-1$. The \emph{size}\index{size} of a \omterm{def:g12:trees:tree}{tree} is its number of \omterm{def:g12:trees:tree}{nodes}.
\end{definition}

\begin{omfigure}
\begin{tikzpicture}[font=\footnotesize, nd/.style={draw, circle, minimum size=6mm, inner sep=0pt, fill=omDef!10},
    lf/.style={nd, fill=omMeth!15}, level distance=11mm, level 1/.style={sibling distance=30mm},
    level 2/.style={sibling distance=15mm}]
  \node[nd] (A) {A}
    child {node[nd] (B) {B} child {node[lf] (D) {D}} child {node[nd] (E) {E} child {node[lf] (G) {G}}
      child[missing]}}
    child {node[nd] (C) {C} child[missing] child {node[lf] (F) {F}}};
  \draw[dashed, omThm, rounded corners=8pt] ($(B)+(-0.45,0.5)$) -- ($(B)+(0.45,0.5)$) -- ($(E)+(0.55,0.1)$)
    -- ($(G)+(0.45,-0.45)$) -- ($(G)+(-0.45,-0.45)$) -- ($(D)+(-0.5,-0.4)$) -- ($(D)+(-0.5,0.3)$) -- cycle;
  \node[omThm, font=\scriptsize] at ($(D)+(-1.3,0.6)$) {subtree B};
  \foreach \d/\y in {0/0, 1/-1.1, 2/-2.2, 3/-3.3} \node[font=\scriptsize, black!60] at (4.2,\y) {depth $\d$};
  \node[font=\scriptsize, right] at ($(A)+(0.4,0)$) {root};
\end{tikzpicture}

{\small A \omterm{def:g12:trees:tree}{tree} of \omterm{def:g12:trees:height}{size} $7$ and \omterm{def:g12:trees:height}{height} $3$. A is the \omterm{def:g12:trees:tree}{root}; D, G and F are the leaves (shaded); B is the \omterm{def:g12:trees:tree}{parent} of D and
E. The dashed line surrounds the \omterm{def:g12:trees:tree}{subtree} of \omterm{def:g12:trees:tree}{root} B: \omterm{def:g12:trees:height}{size} $4$, \omterm{def:g12:trees:height}{height} $2$.}
\end{omfigure}

\section{Binary trees}\label{sec:g12:trees:binary}

\begin{definition}[Binary tree]\label{def:g12:trees:binary-tree}
A \emph{binary tree}\index{binary tree} is either empty, or a \omterm{def:g12:trees:tree}{node} holding a value and two binary \omterm{def:g12:trees:tree}{trees}, its left
\omterm{def:g12:trees:tree}{subtree} and its right \omterm{def:g12:trees:tree}{subtree}.
\end{definition}

The definition is recursive, and so are the functions on \omterm{def:g12:trees:binary-tree}{binary trees}: each one says what to do with the empty \omterm{def:g12:trees:tree}{tree},
and how to combine the answers of the two \omterm{def:g12:trees:tree}{subtrees}. In Python the empty \omterm{def:g12:trees:tree}{tree} is \code{None}.

\omcode{code/grade-12/04-trees/listings/node.py}{1}{28}{\omterm{def:g12:trees:binary-tree}{Binary trees}: a class of \omterm{def:g12:trees:tree}{nodes}, and three recursive
measures.}\label{lst:g12:trees:node}

\omcode{code/grade-12/04-trees/listings/size_height.py}{1}{13}{The \omterm{def:g12:trees:tree}{tree} of the figure.}\label{lst:g12:trees:size}

\omoutput{code/grade-12/04-trees/outputs/size_height.txt}{The numbers read on the figure, and the two conventions for
small \omterm{def:g12:trees:tree}{trees}.}\label{out:g12:trees:size}

\begin{theorem}[Size and height]\label{thm:g12:trees:bounds}
A \omterm{def:g12:trees:binary-tree}{binary tree} of \omterm{def:g12:trees:height}{height} $h \ge 0$ has a \omterm{def:g12:trees:height}{size} $n$ with $h + 1 \le n \le 2^{h+1} - 1$.
\end{theorem}

The lower bound holds because a path from the \omterm{def:g12:trees:tree}{root} to the deepest \omterm{def:g12:trees:tree}{node} already holds $h + 1$ \omterm{def:g12:trees:tree}{nodes}. For the upper
bound, \omterm{def:g12:trees:height}{depth} $d$ holds at most $2^d$ \omterm{def:g12:trees:tree}{nodes}, since each \omterm{def:g12:trees:tree}{node} has at most two children, and
$1 + 2 + 4 + \dots + 2^h = 2^{h+1} - 1$. Exercise \ref{exo:g12:trees:12} proves it again, by induction on the
recursive definition. Turned around: a \omterm{def:g12:trees:binary-tree}{binary tree} of $n$ \omterm{def:g12:trees:tree}{nodes} has a \omterm{def:g12:trees:height}{height} of at least the number of times $n + 1$
can be halved, minus one; $1023$ \omterm{def:g12:trees:tree}{nodes} need a \omterm{def:g12:trees:height}{height} of at least $9$.

\section{Traversals}\label{sec:g12:trees:traversals}

\begin{definition}[Traversal, preorder, inorder, postorder, breadth-first order]\label{def:g12:trees:traversal}
A \emph{traversal}\index{traversal} of a \omterm{def:g12:trees:tree}{tree} visits each \omterm{def:g12:trees:tree}{node} once. In \emph{preorder}\index{preorder}, a \omterm{def:g12:trees:tree}{node} is
visited before its left and right \omterm{def:g12:trees:tree}{subtrees}; in \emph{inorder}\index{inorder}, between them; in
\emph{postorder}\index{postorder}, after them. In \emph{breadth-first order}\index{breadth-first order}, \omterm{def:g12:trees:tree}{nodes} are
visited level by level, from the \omterm{def:g12:trees:tree}{root} down, from left to right in each level.
\end{definition}

\omcode{code/grade-12/04-trees/listings/traversals.py}{1}{32}{Four \omterm{def:g12:trees:traversal}{traversals}; the first three recursive, the last with
a queue.}\label{lst:g12:trees:traversals}

\omoutput{code/grade-12/04-trees/outputs/traversals.txt}{The seven \omterm{def:g12:trees:tree}{nodes} of the figure, in four
orders.}\label{out:g12:trees:traversals}

\omterm{def:g12:trees:traversal}{Breadth-first order} cannot be written as ``a \omterm{def:g12:trees:tree}{node}, then its \omterm{def:g12:trees:tree}{subtrees}'': the children of B must wait while C is
visited. A queue holds the \omterm{def:g12:trees:tree}{nodes} found but not yet visited, and gives them back in the order they were found, level
after level.

An arithmetic expression is a \omterm{def:g12:trees:tree}{tree}: each operator is a \omterm{def:g12:trees:tree}{node} whose \omterm{def:g12:trees:tree}{subtrees} are its operands, and the numbers are the
leaves. Its value is computed in \omterm{def:g12:trees:traversal}{postorder}, operands before operator; \omterm{def:g12:trees:traversal}{postorder} writes the \omterm{def:g12:lists-stacks-queues:postfix}{postfix notation} of the
last chapter, and \omterm{def:g12:trees:traversal}{inorder}, with brackets, the usual one.

\begin{omfigure}
\begin{tikzpicture}[font=\footnotesize, nd/.style={draw, circle, minimum size=6mm, inner sep=0pt, fill=omDef!10},
    lf/.style={nd, fill=omMeth!15}, level distance=11mm, level 1/.style={sibling distance=30mm},
    level 2/.style={sibling distance=14mm}]
  \node[nd] (m) {$\times$}
    child {node[nd] (p) {$+$} child {node[lf] {$3$}} child {node[lf] {$4$}}}
    child {node[nd] (s) {$-$} child {node[lf] {$5$}} child {node[lf] {$2$}}};
  \node[omThm, right, font=\scriptsize] at ($(m)+(0.35,0)$) {$21$};
  \node[omThm, left, font=\scriptsize] at ($(p)+(-0.35,0)$) {$7$};
  \node[omThm, right, font=\scriptsize] at ($(s)+(0.35,0)$) {$3$};
  \node[align=left, font=\scriptsize, right] at (3.3,-1.2) {preorder: $\times$ $+$ $3$ $4$ $-$ $5$ $2$\\
    inorder: $3$ $+$ $4$ $\times$ $5$ $-$ $2$\\ postorder: $3$ $4$ $+$ $5$ $2$ $-$ $\times$\\
    breadth-first: $\times$ $+$ $-$ $3$ $4$ $5$ $2$};
\end{tikzpicture}

{\small The expression $(3 + 4) \times (5 - 2)$ as a \omterm{def:g12:trees:tree}{tree}, with its four \omterm{def:g12:trees:traversal}{traversals}. In \omterm{def:g12:trees:traversal}{postorder} each operator comes
after its operands; the values computed at each \omterm{def:g12:trees:tree}{node} are written beside it.}
\end{omfigure}

\omcode{code/grade-12/04-trees/listings/expression_tree.py}{1}{26}{Evaluating and writing an expression
\omterm{def:g12:trees:tree}{tree}.}\label{lst:g12:trees:expression}

\omoutput{code/grade-12/04-trees/outputs/expression_tree.txt}{The value, and the postfix form read in
\omterm{def:g12:trees:traversal}{postorder}.}\label{out:g12:trees:expression}

\section{Binary search trees}\label{sec:g12:trees:bst}

\begin{definition}[Binary search tree]\label{def:g12:trees:bst}
A \emph{binary search tree}\index{binary search tree} is a \omterm{def:g12:trees:binary-tree}{binary tree} of keys in which, at every \omterm{def:g12:trees:tree}{node}, the keys of
the left \omterm{def:g12:trees:tree}{subtree} are smaller than the \omterm{def:g12:trees:tree}{node}'s key and the keys of the right \omterm{def:g12:trees:tree}{subtree} are larger.
\end{definition}

To find a key, compare it with the \omterm{def:g12:trees:tree}{root}: equal, found; smaller, go on in the left \omterm{def:g12:trees:tree}{subtree}; larger, in the right. To
insert a key, search for it; where the search falls off the \omterm{def:g12:trees:tree}{tree}, hang a new \omterm{def:g12:trees:tree}{leaf}. Both follow one path from the \omterm{def:g12:trees:tree}{root},
so they visit at most $h + 1$ \omterm{def:g12:trees:tree}{nodes} in a \omterm{def:g12:trees:tree}{tree} of \omterm{def:g12:trees:height}{height} $h$. The functions below loop instead of recursing, for a
reason the end of the chapter shows.

\omcode{code/grade-12/04-trees/listings/bst.py}{1}{37}{Insertion and search in a \omterm{def:g12:trees:bst}{binary search
tree}.}\label{lst:g12:trees:bst}

The town library's first ten titles, inserted in the order of its table:

\omcode{code/grade-12/04-trees/listings/bst_demo.py}{1}{17}{A search \omterm{def:g12:trees:tree}{tree} of book titles.}\label{lst:g12:trees:demo}

\omoutput{code/grade-12/04-trees/outputs/bst_demo.txt}{Titles compare alphabetically; the \omterm{def:g12:trees:traversal}{inorder} \omterm{def:g12:trees:traversal}{traversal} lists them
sorted.}\label{out:g12:trees:demo}

\begin{omfigure}
\begin{tikzpicture}[font=\scriptsize, nd/.style={draw, rounded corners=2pt, fill=omDef!10, inner sep=2pt,
    align=center}, level distance=10mm, level 1/.style={sibling distance=50mm},
    level 2/.style={sibling distance=40mm}, level 3/.style={sibling distance=44mm},
    level 4/.style={sibling distance=26mm}]
  \node[nd] {The Silent Code}
    child {node[nd] {The Paper Forest}
      child {node[nd] {A Year by the Island}
        child {node[nd] {A Map of\\ the Lighthouse} child[missing] child {node[nd] {A Year by\\ the Bridge}}}
        child {node[nd] {Seven Library} child {node[nd] {Seven Code}}
          child {node[nd] {The Lost Forest} child {node[nd, fill=omThm!12] {The Clockwork\\ Harbour}}
            child[missing]}}}
      child[missing]}
    child {node[nd] {The Silent Kingdom}};
\end{tikzpicture}

{\small The search \omterm{def:g12:trees:tree}{tree} of \code{bst_demo.py}. Every title on the left of a \omterm{def:g12:trees:tree}{node} comes before it in alphabetical
order, every title on the right after it. The deepest \omterm{def:g12:trees:tree}{node}, shaded, has \omterm{def:g12:trees:height}{depth} $5$: the \omterm{def:g12:trees:height}{height} of the \omterm{def:g12:trees:tree}{tree}.}
\end{omfigure}

\begin{remark}
Alphabetical order here is Python's order on strings: character by character, by \omterm{def:g10:text-encodings:code-point}{code point}. Every capital letter comes
before every small one: \code{"Zoe"} comes before \code{"apple"}.
\end{remark}

\section{Balanced and degenerate trees}\label{sec:g12:trees:balance}

\begin{definition}[Balanced tree, degenerate tree]\label{def:g12:trees:balanced}
A \omterm{def:g12:trees:binary-tree}{binary tree} is \emph{balanced}\index{balanced tree} when its \omterm{def:g12:trees:height}{height} is close to the smallest possible for its \omterm{def:g12:trees:height}{size},
about the \omterm{def:g11:cost-of-an-algorithm:halvings}{number of halvings} of $n$. It is \emph{degenerate}\index{degenerate tree} when every \omterm{def:g12:trees:tree}{node} has at most one
\omterm{def:g12:trees:tree}{child}: it is then a list, of \omterm{def:g12:trees:height}{height} $n - 1$.
\end{definition}

The shape of a search \omterm{def:g12:trees:tree}{tree} depends only on the order of insertion. Keys inserted in increasing order each go to the
right of all the others: a \omterm{def:g12:trees:balanced}{degenerate tree}.

\begin{omfigure}
\begin{tikzpicture}[font=\footnotesize, nd/.style={draw, circle, minimum size=5.5mm, inner sep=0pt, fill=omDef!10},
    level distance=8mm]
  \begin{scope}[level 1/.style={sibling distance=7mm}]
  \node[nd] {1} child[missing] child {node[nd] {2} child[missing] child {node[nd] {3} child[missing]
    child {node[nd] {4} child[missing] child {node[nd] {5} child[missing] child {node[nd] {6} child[missing]
    child {node[nd] {7}}}}}}};
  \node[font=\scriptsize] at (2.4,-5.4) {inserted $1, 2, \dots, 7$: height $6$};
  \end{scope}
  \begin{scope}[xshift=8cm, level 1/.style={sibling distance=24mm}, level 2/.style={sibling distance=12mm},
      level distance=11mm]
  \node[nd] {4} child {node[nd] {2} child {node[nd] {1}} child {node[nd] {3}}}
    child {node[nd] {6} child {node[nd] {5}} child {node[nd] {7}}};
  \node[font=\scriptsize] at (0,-3.2) {inserted $4, 2, 6, 1, 3, 5, 7$: height $2$};
  \end{scope}
\end{tikzpicture}

{\small The same seven keys, two orders of insertion. Finding $7$ visits seven \omterm{def:g12:trees:tree}{nodes} on the left, three on the right.}
\end{omfigure}

\omcode{code/grade-12/04-trees/listings/degenerate.py}{1}{20}{\omterm{def:g12:trees:height}{Heights} after three orders of seven keys, then $1023$
keys sorted and shuffled.}\label{lst:g12:trees:degenerate}

\omoutput{code/grade-12/04-trees/outputs/degenerate.txt}{The sorted \omterm{def:g12:trees:tree}{tree} of $1023$ keys is so deep that the recursive
\code{height} fails.}\label{out:g12:trees:degenerate}

The recursive \code{height} calls itself once per level: on a list of $1023$ \omterm{def:g12:trees:tree}{nodes}, that is beyond Python's limit of
$1000$ frames (chapter \ref{ch:g12:recursion}). This is why \code{insert} and \code{search} loop. \omterm{def:g12:trees:tree}{Trees} that keep
themselves \omterm{def:g12:trees:balanced}{balanced} whatever the order, by reshaping themselves at each insertion, exist: they guarantee a \omterm{def:g12:trees:height}{height} in
proportion to the \omterm{def:g11:cost-of-an-algorithm:halvings}{number of halvings} of $n$, and belong to a later volume.

\section{Exercises}\label{sec:g12:trees:exercises}

\begin{exercise}[$\star$]\label{exo:g12:trees:1}
Give the \omterm{def:g12:trees:height}{size}, \omterm{def:g12:trees:height}{height} and number of leaves of
\code{Node(1, Node(2, Node(4), Node(5, Node(7), Node(8))), Node(3, Node(6)))}.
\end{exercise}

\begin{exercise}[$\star$]\label{exo:g12:trees:2}
Give the \omterm{def:g12:trees:traversal}{preorder} of the \omterm{def:g12:trees:tree}{tree} of exercise \ref{exo:g12:trees:1}.
\end{exercise}

\begin{exercise}[$\star$]\label{exo:g12:trees:3}
Give its \omterm{def:g12:trees:traversal}{inorder} and its \omterm{def:g12:trees:traversal}{postorder}.
\end{exercise}

\begin{exercise}[$\star$]\label{exo:g12:trees:4}
Draw the search \omterm{def:g12:trees:tree}{tree} obtained by inserting $50$, $30$, $70$, $20$, $40$, $60$, $80$, $35$, in that order. What is its
\omterm{def:g12:trees:height}{height}?
\end{exercise}

\begin{exercise}[$\star$]\label{exo:g12:trees:5}
Write \code{count_at_depth(t, d)}, the number of \omterm{def:g12:trees:tree}{nodes} of \omterm{def:g12:trees:height}{depth} \code{d} in the \omterm{def:g12:trees:tree}{tree} \code{t}.
\end{exercise}

\begin{exercise}[$\star\star$]\label{exo:g12:trees:6}
A student writes \code{height} with \code{return 0} for the empty \omterm{def:g12:trees:tree}{tree}. What does it return for a single \omterm{def:g12:trees:tree}{node}, and
for the \omterm{def:g12:trees:tree}{tree} of the first figure? How must the bounds of the theorem be rewritten for that function?
\end{exercise}

\begin{exercise}[$\star\star$]\label{exo:g12:trees:7}
Write \code{bst_min(t)}, the smallest key of a non-empty search \omterm{def:g12:trees:tree}{tree}, without visiting more than one path.
\end{exercise}

\begin{exercise}[$\star\star$]\label{exo:g12:trees:8}
What is the smallest possible \omterm{def:g12:trees:height}{height} of a \omterm{def:g12:trees:binary-tree}{binary tree} of $100$ \omterm{def:g12:trees:tree}{nodes}?
\end{exercise}

\begin{exercise}[$\star\star$]\label{exo:g12:trees:9}
Give the \omterm{def:g12:trees:traversal}{breadth-first order} of the \omterm{def:g12:trees:tree}{tree} of exercise \ref{exo:g12:trees:4}.
\end{exercise}

\begin{exercise}[$\star\star$]\label{exo:g12:trees:10}
This recursive insertion is called as \code{t = insert(t, k)}. Inserting $50$, $30$, $70$ leaves a \omterm{def:g12:trees:tree}{tree} of one \omterm{def:g12:trees:tree}{node}.
Why? Correct it.

\smallskip\noindent\hspace*{2em}\code{def insert(t, key):}

\noindent\hspace*{4em}\code{if t is None:}

\noindent\hspace*{6em}\code{return Node(key)}

\noindent\hspace*{4em}\code{if key < t.value:}

\noindent\hspace*{6em}\code{insert(t.left, key)}

\noindent\hspace*{4em}\code{else:}

\noindent\hspace*{6em}\code{insert(t.right, key)}

\noindent\hspace*{4em}\code{return t}
\end{exercise}

\begin{exercise}[$\star\star$]\label{exo:g12:trees:11}
Write \code{is_bst(t)}. Show that checking, at each \omterm{def:g12:trees:tree}{node}, only that the left \omterm{def:g12:trees:tree}{child} is smaller and the right \omterm{def:g12:trees:tree}{child}
larger is not enough.
\end{exercise}

\begin{exercise}[$\star\star\star$]\label{exo:g12:trees:12}
Prove, by induction on the recursive definition, that a \omterm{def:g12:trees:binary-tree}{binary tree} of \omterm{def:g12:trees:height}{height} $h$ and \omterm{def:g12:trees:height}{size} $n$ satisfies
$n \le 2^{h+1} - 1$.
\end{exercise}

\begin{exercise}[$\star\star\star$]\label{exo:g12:trees:13}
Prove that the \omterm{def:g12:trees:traversal}{inorder} \omterm{def:g12:trees:traversal}{traversal} of a \omterm{def:g12:trees:bst}{binary search tree} lists its keys in increasing order.
\end{exercise}

\begin{exercise}[$\star\star\star$]\label{exo:g12:trees:14}
A \omterm{def:g12:trees:tree}{tree} has distinct values. Write \code{rebuild(pre, ino)}, which rebuilds it from its \omterm{def:g12:trees:traversal}{preorder} and its \omterm{def:g12:trees:traversal}{inorder}.
\end{exercise}

\begin{exercise}[$\star\star\star$]\label{exo:g12:trees:15}
How many \omterm{def:g12:trees:binary-tree}{binary trees} of \omterm{def:g12:trees:height}{size} $3$ are there, counting shapes only? Of \omterm{def:g12:trees:height}{size} $4$? Write a \omterm{def:g12:recursion:recursion}{recursive function} that counts
those of \omterm{def:g12:trees:height}{size} $n$.
\end{exercise}

\section{Problem: How Tall Is a Search Tree?}\label{sec:g12:trees:problem}

\begin{problem}[{Weekend problem --- the height of a search tree}]\label{pb:g12:trees:1}
The keys $1, 2, \dots, 1023$ are inserted into a search \omterm{def:g12:trees:tree}{tree}, in an order that decides its \omterm{def:g12:trees:height}{height}. The program below
builds $1000$ \omterm{def:g12:trees:tree}{trees} from $1000$ random orders of the same keys, and counts the \omterm{def:g12:trees:height}{heights}; the chart draws the counts.

\omcode{code/grade-12/04-trees/listings/random_heights.py}{1}{23}{\omterm{def:g12:trees:height}{Heights} of $1000$ random search
\omterm{def:g12:trees:tree}{trees}.}\label{lst:g12:trees:random}

\omoutput{code/grade-12/04-trees/outputs/random_heights.txt}{The smallest, largest and mean \omterm{def:g12:trees:height}{heights}, then how many
\omterm{def:g12:trees:tree}{trees} had each \omterm{def:g12:trees:height}{height}.}\label{out:g12:trees:random}

\begin{omfigure}
\begin{tikzpicture}
\begin{axis}[omrecord, ybar, width=12cm, height=5cm, bar width=4mm, xlabel={height}, ylabel={trees}, ymin=0,
    xtick={16,18,...,30}]
  \addplot[fill=omDef!40, draw=omDef] table[x=height, y=count]{code/grade-12/04-trees/figures/heights_hist.dat};
\end{axis}
\end{tikzpicture}

{\small \omterm{def:g12:trees:height}{Heights} of the $1000$ search \omterm{def:g12:trees:tree}{trees} of $1023$ keys, from \code{figures/heights_hist.py}. None comes near the
worst case of $1022$; none reaches the best case of $9$ either.}
\end{omfigure}

\medskip\noindent\textbf{Part I --- A search tree of $n$ keys.}
\begin{enumerate}
  \item Why does finding a key in a search \omterm{def:g12:trees:tree}{tree} of \omterm{def:g12:trees:height}{height} $h$ visit at most $h + 1$ \omterm{def:g12:trees:tree}{nodes}?
  \item Does the same hold for a key that is absent?
  \item Why is the \omterm{def:g12:trees:traversal}{inorder} \omterm{def:g12:trees:traversal}{traversal} of every one of the $1000$ \omterm{def:g12:trees:tree}{trees} the same?
  \item Which key is at the \omterm{def:g12:trees:tree}{root} of each \omterm{def:g12:trees:tree}{tree}?
\end{enumerate}

\medskip\noindent\textbf{Part II --- The worst case.}
\begin{enumerate}[resume]
  \item Show that inserting $1, 2, \dots, n$ in increasing order gives a \omterm{def:g12:trees:tree}{tree} of \omterm{def:g12:trees:height}{height} $n - 1$.
  \item Give another order of $1, \dots, 7$ that gives \omterm{def:g12:trees:height}{height} $6$, other than increasing and decreasing.
  \item Why does the recursive \code{height} fail on the sorted \omterm{def:g12:trees:tree}{tree} of $1023$ keys, while \code{search} works?
  \item How many \omterm{def:g12:trees:tree}{nodes} does finding $1023$ visit in the sorted \omterm{def:g12:trees:tree}{tree}, and why?
\end{enumerate}

\medskip\noindent\textbf{Part III --- The best case.}
\begin{enumerate}[resume]
  \item From the theorem, what is the smallest \omterm{def:g12:trees:height}{height} of a \omterm{def:g12:trees:tree}{tree} of $1023$ \omterm{def:g12:trees:tree}{nodes}?
  \item Describe an order of insertion of $1, \dots, 1023$ that reaches it.
  \item In that \omterm{def:g12:trees:tree}{tree}, how many \omterm{def:g12:trees:tree}{nodes} does any search visit at most?
  \item Write that order for $1, \dots, 7$, and check it against the output of \code{degenerate.py}.
\end{enumerate}

\medskip\noindent\textbf{Part IV --- Random orders.}
\begin{enumerate}[resume]
  \item Read the smallest, largest and mean \omterm{def:g12:trees:height}{heights} of the $1000$ random \omterm{def:g12:trees:tree}{trees}.
  \item Which \omterm{def:g12:trees:height}{height} is the most frequent? How many \omterm{def:g12:trees:tree}{trees} out of $1000$ had a \omterm{def:g12:trees:height}{height} of $25$ or more?
  \item Compare the mean with the best \omterm{def:g12:trees:height}{height} $9$ and the worst $1022$: how many times larger than the best, and how
    many times smaller than the worst?
  \item Why is a \omterm{def:g12:trees:height}{height} of $1022$ possible, but never seen?
  \item Searching a random \omterm{def:g12:trees:tree}{tree} visits at most how many \omterm{def:g12:trees:tree}{nodes}, in the tallest of the $1000$?
  \item Why is the experiment built from a fixed seed?
  \item The library inserts its titles in the order of its table. Is that order random with respect to the titles?
    What would happen if the table were sorted by title?
  \item State the result: the best \omterm{def:g12:trees:height}{height}, the worst \omterm{def:g12:trees:height}{height}, and the measured \omterm{def:g12:trees:height}{heights}.
\end{enumerate}
\end{problem}
